% ============================================================
\section{Problema 4: VLSM III}
% ============================================================

\begin{tcolorbox}[enunciado]
Debido a la gran expansión en el flujo de pasajeros internacionales, operaciones de carga comercial y la incorporación de distritos hoteleros y comerciales, la administración de TI del \textbf{Aeropuerto Internacional Felipe Ángeles (AIFA)} requiere reestructurar por completo su direccionamiento lógico de red.

El direccionamiento debe realizarse de manera óptima utilizando la técnica de \textbf{VLSM (Variable Length Subnet Mask)} a partir del bloque de direccionamiento privado de Clase A: \textbf{10.0.0.0/8}.

Para garantizar la estabilidad del enrutamiento y evitar el desperdicio de direcciones, se solicita diseñar la topología considerando los requerimientos de hosts de cada infraestructura y servicio del aeropuerto. Los datos deben ser organizados estrictamente de mayor a menor número de hosts antes de asignar los segmentos, incluyendo las interconexiones entre los enrutadores principales de la red troncal (backbone).

A continuación se presenta la distribución oficial calculada para los departamentos, terminales y enlaces troncales del complejo aeroportuario:

\vspace{0.3cm}
\begin{center}
    \begin{tabular}{lc}
    \toprule
    \textbf{Infraestructura (AIFA)} & \textbf{Hosts} \\
    \midrule
    Terminal 1 Pasajeros & 8000 \\
    Red Wifi Pública Principal & 5000 \\
    Terminal 2 Pasajeros & 4100 \\
    Aduanas e Impuestos T1 & 3500 \\
    Red de Locales Comerciales & 3000 \\
    Zona de Carga y Logística & 2600 \\
    Oficinas de Aerolíneas & 2200 \\
    Hoteles y Alojamientos AIFA & 2050 \\
    Sistemas de Check-in T2 & 1800 \\
    Monitoreo de Cámaras CCTV & 1500 \\
    Estacionamientos y Accesos & 1200 \\
    Torre de Control y Pistas & 1000 \\
    Centro de Datos AIFA & 500 \\
    Enlace Core a Terminal 1 & 2 \\
    Enlace Core a Terminal 2 & 2 \\
    Enlace Core a Torre de Control & 2 \\
    Enlace Core a Zona de Carga & 2 \\
    Enlace Terminal 1 a Terminal 2 & 2 \\
    Enlace Respaldo Data Center & 2 \\
    \bottomrule
    \end{tabular}
\end{center}
\vspace{0.3cm}

Se debe hacer el VLSM con esos hosts.
\end{tcolorbox}

\subsection*{Solución}
% Escribe aquí tu solución para el Problema 4...